\section{Objetivos}
\label{sc:objetivos}

\subsection{Objetivo Geral}
Desenvolver, implementar e validar uma arquitetura de inteligência geoespacial ativa (\textit{Actionable GeoAI DataOps}) para mitigação do risco de incêndios sobre infraestruturas do SIN e áreas protegidas, integrando modelagem anisotrópica por vento, indexação Uber H3, alarmes via Webhooks e agentes de IA para relatórios situacionais.

\subsection{Objetivos Específicos}
\begin{enumerate}
    \setlength\itemsep{0em}
    \item Integrar dados horários de vento (ECMWF ERA5 / INMET) e histórico de perturbações do ONS ao pipeline Medallion;
    \item Discretizar feições e calcular o risco espacial via células hexagonais Uber H3 (\texttt{uint64}), substituindo junções geométricas por \textit{hash joins} $O(1)$;
    \item Desenvolver e calibrar o algoritmo de buffers elípticos anisotrópicos governado pelo modelo de Rothermel e classes de combustível do MapBiomas Fogo;
    \item Construir o motor de alarmes com máquina de estados de \textit{cooldown} (anti-spam) e despacho assíncrono via Webhooks (Telegram e HTTP/REST);
    \item Implementar servidor FastMCP dotado de ferramentas determinísticas (\textit{Text-to-SQL} no Apache Druid e grafo elétrico NetworkX) para geração automática de relatórios executivos no padrão ICS-209 com garantia de zero alucinação;
    \item Avaliar a acurácia espacial, a latência de triagem ($< 30\text{ s}$) e a fidelidade factual do sistema frente a métodos isotrópicos tradicionais.
\end{enumerate}
